\chapter{Related Research and Analytic Concepts}\label{chap:2}

Shared-view interfaces change both the information available to collaborators and the ways they can act on it. This chapter connects research on visual access with accounts of how people develop and revise a shared task. The comparison centres on three questions: what a view makes inspectable, how a partner's actions become usable, and how contributions remain connected as control changes.

\section{Seeing the workspace and understanding a contribution}\label{sec:2-1}

A common image gives partners access to a task state, while their responses establish what they understand about it. Clark and Brennan describe grounding as establishing that a contribution has been understood sufficiently for the current purpose \cite{clark1991}. Gergle, Kraut and Fussell distinguish this work from maintaining awareness of the current situation; their experiments examine visual immediacy, perspective alignment and common viewing area \cite{gergle2013}. The distinction is useful in construction tasks. Seeing where a piece has been placed, understanding the proposed relation and accepting the whole arrangement as correct are separate achievements.

The shared state also changes as people work. Roschelle and Teasley's joint problem space connects goals, descriptions of the problem and available actions through the partners' ongoing activity \cite{roschelle1995}. Gutwin and Greenberg explain how bodies, communication and the visible effects of manipulation provide workspace awareness. Their concept of feedthrough identifies the information that object changes make available to others \cite{gutwin2002}. Read together, these accounts connect a placement to the partner's next response. An arrangement can communicate an idea, provide a reference for another move or become the part that the pair later reconsiders.

Recent studies make the changing demand for this information visible. In six VR--desktop pairs, Saffo et al. found that sharing interactions and perspectives helped partners follow one another when moving towards closer collaboration \cite{saffo2023}. Irlitti et al.'s six three-person groups in hybrid workspaces showed that one person's manipulation could accompany either active discussion by others or relatively disengaged watching \cite{irlitti2024}. Both studies connect participation to the surrounding activity. For Tangram, this supports examining what a move responds to and what it enables next, alongside who holds the piece.

\section{What different shared views support}\label{sec:2-2}

Shared perspective is implemented through several distinct mechanisms. TeamPortal provides an interactive parallel view into a partner's field of view \cite{wang2025}. MAGIC transforms the partner's representation so that gestures remain meaningful across spatial arrangements while retaining face-to-face interaction \cite{fidalgo2023}. SPARC combines a common orientation with transformed avatar arms and references; its evaluation uses assigned assembler and instructor roles in three-dimensional assembly \cite{simoes2025}. These systems respectively emphasise access to another view, alignment of embodied references and persistent shared orientation. They show why camera position, partner representation and control mapping need separate description when comparing interfaces.

The value of an overview depends on the spatial demands of the task. Chen et al. compared a world-in-miniature, an interactive 2D map and a static terrain map with 18 pairs. Interactive views improved performance over the baseline, and the miniature was more useful in complex scenes involving occlusion \cite{chen2024}. Chang, Lee and Yoo compared 2D video, a 360-degree view and a reconstructed 3D view in VR--AR assistance. Thirty participants each worked with an actor; the independently explorable 3D view supported faster guidance, while the 360-degree view was preferred for the feeling of being together \cite{chang2024}. Both comparisons link richer access to spatial information with particular task benefits, while Chang et al. also separate those benefits from social experience.

These comparisons involve different forms of control and cooperation. Chen et al.'s map supports horizontal movement, whereas the miniature supports three-dimensional operations. Chang et al.'s task assigns expert and worker roles. In a constrained passenger-space assembly task, Medeiros et al. instead varied the users' spatial arrangement. Side-by-side placement was preferred for visibility and object control, and face-to-face placement offered stronger social presence; completion time and correct-piece counts did not differ significantly \cite{medeiros2024}. A favourable experience of the view therefore needs to be related to the action and outcome it supports.

Communication frequency requires similar interpretation. Basille, Lavou\'e and Serna's desktop serious-game study with six triads associated first-person viewing with more information sharing, which the authors related to differences in what players could see \cite{basille2022}. This provides a useful comparison with shared-overview designs: additional communication can supply information missing from a participant's view. Its frequency alone says little about whether the pair can judge the developing configuration accurately.

Tangram makes precise planar relationships central to this comparison. Chen et al. identify overlap in a 2D projection as a problem, and their placement task leaves finer placement accuracy for further investigation \cite{chen2024}. In Tangram, relative widths, orientations and edges determine whether a candidate matches the target. A perpendicular orthographic view preserves these planar relations without perspective depth scaling. The present study examines how partners use the available view while forming, inspecting and accepting a construction, bringing geometric judgment into the analysis of collaborative action.

\section{Connecting visual access with control}\label{sec:2-3}

Partners can understand the same construction while having different practical access to its pieces. Irlitti et al.'s physical, volumetric and virtual objects offered different possibilities for direct action \cite{irlitti2024}. Medeiros et al.'s spatial formations also changed visibility and object control \cite{medeiros2024}. These studies connect the organisation of participation to the way the interface distributes action. This is relevant to the present overhead prototype, where the shared camera image and the tracked movement of each controller use different frames of reference. A clearly visible next step may still be difficult for one partner to execute.

Making the current action state visible can help coordinate that execution. In a screen-based virtual environment with hand tracking, Kim and Park examined feedback about the user's and partner's states during joint movement and handover \cite{kim2024}. Their study addresses the timing of coordinated manipulation. Combined with the changing awareness needs found by Saffo et al. \cite{saffo2023}, it suggests examining assistance as an episode: one participant encounters a difficulty, a partner gains enough information to act, and control or responsibility changes for the next step. The relevant question is how that transfer advances the shared construction.

\section{Developing ideas through temporary divisions of work}\label{sec:2-4}

Manipulation can help develop a proposal as well as communicate it. Hornecker and Buur connect lightweight, legible manipulation with the social organisation of tangible interaction \cite{hornecker2006}. Kirsh and Maglio's analysis of Tetris explains how moving or rotating a piece can make a spatial problem easier to assess \cite{kirsh1994}. Their accounts give complementary reasons to examine tentative placements: the action can help its maker work out a relation and make that relation available to a partner. Kirsh and Maglio also recognise that the same move may advance the task and help determine what to do. Its role is interpreted through the surrounding sequence.

The resulting division of work can be brief. Dillenbourg discusses cooperation and collaboration in terms of the grain and organisation of activity, including spontaneous division and changing, intertwined roles \cite{dillenbourg1999}. Gutwin and Greenberg similarly address movement between independent activity and closer joint work \cite{gutwin2002}. These accounts accommodate a pair that develops different parts for a time and then returns to a common difficulty. The connection between contributions lies in how one person's work extends, adjusts or replaces what the other has made.

Measurement needs to preserve this distinction. In Yang et al.'s VR--desktop image-organisation study, 40 participants worked in groups of four. VR groups showed more conversation and object interaction, with more equal contributions to grabbing, while speed and accuracy were similar \cite{yang2022}. Their measures distinguish picking up an image from regrouping it. Irlitti et al.'s distinction between engaged discussion and watching further shows why a count of direct operations needs contextual interpretation \cite{irlitti2024}. For Tangram, an accepted grab identifies control; the resulting configuration and the partner's response establish whether it was a small adjustment, a new proposal or a structural revision.

\section{Keeping and revising a partner's work}\label{sec:2-5}

A construction remains available after its builder releases it. Preserving, extending and replacing that arrangement have different consequences for the shared task. Yang et al. examined requests for permission before changing another person's image grouping \cite{yang2022}. Rasch et al. studied undo with 16 pairs building towers in VR: world undo increased reported connectedness while also creating opportunities to interfere with a partner's work \cite{rasch2024}. Together, these studies make both the surrounding exchange and the extent of a change relevant. Retained geometry establishes that earlier work remains available; talk or retrospective accounts can explain what that retention means to the participants.

Rasch et al. also identify changing collaboration arrangements as an issue beyond their deliberately divided and collaborative task settings \cite{rasch2024}. Dillenbourg's account of intertwined work helps connect this issue to short construction episodes \cite{dillenbourg1999}. One person may build a part, another may continue it, and the original builder may later reopen it. A local alternative can redirect the next step before either partner has a complete solution. Following who formed and changed the arrangement makes those developments visible.

This thesis connects these questions through the evolving Tangram configuration. Viewpoint research establishes why task geometry, orientation and control matter. Grounding and workspace awareness explain why the response to a placement matters. Research on manipulation, division of work and revision provides concepts for examining how partial contributions persist and change. The study follows these relationships across physical and VR conditions to identify what shared-overhead VR needs to support.
